\section{Marco Teórico}

\subsection{Retinopatía Diabética e Importancia Médica de las Lesiones}

La retinopatía diabética (RD) es una de las principales causas de ceguera en el mundo y se caracteriza por ser una ``enfermedad silenciosa'' que progresa sin síntomas hasta etapas críticas \parencite{casanova2014application}. 

Ya en el año 2003 el estudio de la retinopatía diabética hacía patente que esta enfermedad no se diagnostica por cambios de nervio óptico, sino que por una serie de lesiones jerárquicas: primero aparecen los microaneurismas, luego las hemorragias y por último los exudados. \parencite{wilkinson2003proposed}. En este punto entendemos la riqueza en la información proporcionada por el dataset utilizado. 

Importante recalcar que gracias al procesamiento digital de imágenes retinales y a los modelos de aprendizaje automático es que podemos llegar a saber qué partes del ojo delatan esta enfermedad y también qué tan avanzada está.

Existen dos indicadores clave para clasificar a un paciente como sano o enfermo: los microaneurismas y los exudados.

\begin{itemize}
    \item \textbf{Microaneurismas (MA)}: Son las lesiones clínicas más tempranas de la RD. Consisten en pequeñas dilataciones en las paredes de los capilares retinales, causadas por el debilitamiento estructural provocado por niveles elevados de glucosa en la sangre. En imágenes de fondo de ojo (retinografía) se visualizan como puntos rojos diminutos y circulares con bordes bien definidos, es interesante indicar que el  contraste frente al fondo anaranjado de la retina permite su identificación mediante filtros morfológicos \parencite{antal2014ensemble}. Su conteo es el principal indicador de la etapa inicial de la enfermedad y el factor más relevante para discriminar computacionalmente la presencia de RD \parencite{casanova2014application}.

    \item \textbf{Exudados}: Depósitos que aparecen cuando los vasos dañados filtran fluidos hacia el tejido retinal, indicando avance de la enfermedad hacia etapas moderadas o severas. En el análisis de datos, el área de los exudados se normaliza según el diámetro de la Región de Interés (ROI) para compensar variaciones en el tamaño de las imágenes \parencite{antal2014dataset}. Podemos encontrar dos tipos:
    \begin{itemize}
        \item \textbf{Exudados duros}: Son depósitos de lípidos (grasas) y proteínas que se filtran desde el torrente sanguíneo hacia la retina. Se presentan como manchas pequeñas, brillantes y de tono amarillento.
        \item \textbf{Exudados blandos (manchas algodonosas)}: Áreas de infarto (isquemia) donde el tejido nervioso muere por falta de oxígeno. Podemos identificarlas como manchas blancas más grandes con bordes difusos, indicativas de un avance en la enfermedad más severo.
    \end{itemize}
    \textit{Nota metodológica:} La distinción clínica entre exudados duros y blandos se incluye en este trabajo con fines estrictamente informativos. Para efectos del análisis computacional y el procesamiento del dataset, se utiliza el bloque consolidado de variables que comprende desde \texttt{exudate1} hasta \texttt{exudate8}, las cuales capturan la densidad total de estas lesiones sin realizar una separación explícita en la clasificación del conjunto de datos.
\end{itemize}

\subsection{Fundamentos de Estadística Descriptiva y Forma de Distribución}

En estadística descriptiva buscamos indicadores que sean claves para comprender el comportamiento de las variables del conjunto de datos. En este análisis se responden preguntas como ¿dónde están concentrados los datos? ¿qué tan dispersos están? ¿hay valores extremos?

\textit{``Este análisis descriptivo permitirá confirmar si la distribución de los microaneurismas y exudados presenta una variabilidad significativa entre pacientes sanos y enfermos, validando su importancia como indicadores diagnósticos''}

\begin{itemize}
    \item \textbf{Medidas de tendencia central y dispersión}: Parámetros básicos que nos permiten comprender qué tan variados son los datos de nuestro conjunto, qué tan diferentes son entre sí; este análisis nos permite descubrir patrones de enfermedad. 
    
    \textbf{Media aritmética}: Es el promedio tradicional. Nos da una idea del valor \enquote{típico} de la variable sumando todos los casos y dividiendo por el total.

    \textbf{Mediana}: Es el valor central cuando ordenamos los datos de menor a mayor. A diferencia de la media, es una medida robusta, lo que significa que no se ve afectada por valores extremadamente altos o bajos que no permitirian representar el comportamiento general de los datos.

    \textbf{Desviación estándar}: Nos indica qué tan diferentes son los datos respecto al promedio. Una desviación alta significa que los datos están muy dispersos, que pueden existir valores atípicos con respecto al promedio; una baja, que están muy agrupados cerca del centro (promedio), valores que se asemejen más. Respondemos a la pregunta: ¿Qué tan diferentes son los pacientes entre sí respecto al promedio?

    \item \textbf{Forma de la distribución}: No basta con saber el centro; hay que saber cómo se ``dibuja'' el conjunto de datos.

    \textbf{Asimetría (\textit{Skewness})}: Nos dice si los datos se inclinan hacia un lado. Nos indica si la mayoría de los pacientes tienen valores bajos o altos de una lesión.  Se responde a la pregunta: ¿Hacia dónde se inclina la balanza de pacientes? (Sanos vs. Graves).

    \textbf{Curtosis}: Esta medición nos sirve para encontrar las sorpresas en el conjunto de datos, como por ejemplo, si hay muchos pacientes con valores similares (alta curtosis) o concentración de datos, inferimos entonces que la mayoría de los pacientes tienen valores de la variable similares o iguales, pero hay una probabilidad real de encontrar casos con valores que salen de lo común y se van a los extremos. Es una medición que nos invita a fijarnos en los datos extremos para que podamos identificarlos y que estos no afecten nuestros cálculos. Podemos encontrar dónde están concentrados nuestros datos y también podemos saber si estos casos extremos son raros o son comunes dependiendo de si la curtosis es alta o baja. 

    \item \textbf{Detección de valores atípicos (\textit{Outliers})}: ¿Qué es un outlier? Es un valor que se aleja tanto del resto que puede parecer no pertenecer al mismo grupo o conjunto de datos. Una vez que la asimetría y la curtosis nos advierten que existen datos extremos, utilizamos el Método de Tukey para identificarlos con precisión.
    
    ¿Por qué importa detectarlos? Porque pueden distorsionar el análisis. Podríamos obtener el promedio de un conjunto de datos y no representar la realidad del mismo al considerar datos que salen de lo común y son poco representativos con respecto al conjunto general.

    ¿Qué es el método de Tukey?

    Es la técnica más usada para detectar outliers de forma objetiva. Funciona de la siguiente manera:

    \begin{enumerate}
        \item Toma el 25\% inferior de los datos (Q1) y el 75\% superior (Q3).
        \item Calcula la distancia entre ellos: $IQR = Q3 - Q1$.
        \item Define un límite inferior: $Q1 - 1.5 \times IQR$.
        \item Define un límite superior: $Q3 + 1.5 \times IQR$.
        \item Todo lo que queda fuera de esos límites es un outlier.
    \end{enumerate}
\end{itemize}

\subsection{Fundamentos de Estadística Inferencial y Pruebas de Hipótesis}

  Mientras que el análisis descriptivo solo describe lo que vemos, la inferencial va más allá — te permite sacar conclusiones y tomar decisiones basadas en los datos. La pregunta clave es: ¿lo 
  que observo en mi muestra es real o es producto de la suerte?                                                                                                                         
   
  Por ejemplo: veo que los pacientes enfermos tienen más microaneurismas que los sanos. Pero ¿esa diferencia es real o simplemente mi elección al seleccionar los pacientes del estudio no fue la mejor?                                                                                                                                                                        
                                                                                                                                                                                    
  Las pruebas de hipótesis responden esa pregunta matemáticamente.

  Entre las pruebas de hipótesis que se utilizan en este análisis se encuentran:

\begin{itemize}
    \item \textbf{Prueba de normalidad (Shapiro-Wilk)}:

    Antes de aplicar cualquier prueba estadística, necesitas saber cómo se distribuyen nuestros datos — si siguen una distribución normal o no.

    ¿Por qué importa? Porque muchas pruebas estadísticas asumen que los datos son normales. Si no lo son, debemos usar pruebas distintas.

    Shapiro-Wilk responde: ¿mis datos se comportan como una distribución normal?

    Para responder la pregunta anterior es que tenemos el p-valor, una medida que nos permite decidir si rechazamos o no la hipótesis de normalidad y su límite es 0.05.

    \begin{itemize}
        \item Si p-valor $< 0.05$ $\rightarrow$ No son normales $\rightarrow$ usas pruebas no paramétricas (Pruebas más robustas).
        \item Si p-valor $> 0.05$ $\rightarrow$ Podrían ser normales $\rightarrow$ puedes usar pruebas paramétricas.
    \end{itemize}


    \item \textbf{Comparación de grupos (prueba Mann-Whitney U)}:

    Una vez que sabemos que los datos no son normales, necesitamos una prueba que compare dos grupos sin asumir normalidad. Es la herramienta que decide si hay una diferencia real entre los sanos y los enfermos.

    Mann-Whitney U responde: ¿la distribución de esta variable es significativamente distinta entre pacientes sanos y enfermos?

    \begin{itemize}
        \item Si p-valor $< 0.05$ $\rightarrow$ Sí hay diferencia real entre los grupos.
        \item Si p-valor $> 0.05$ $\rightarrow$ No hay evidencia de diferencia.
    \end{itemize}

    \item \textbf{Homocedasticidad (Test de Levene)}:

    Se utiliza para evaluar la igualdad de varianzas entre diferentes grupos, en este caso grupo de pacientes sanos y enfermos.

    Homocedasticidad es una palabra técnica que significa simplemente: ¿las varianzas de los dos grupos son iguales?
    
    La varianza mide qué tan dispersos están los datos. Si los pacientes sanos tienen todos valores similares de microaneurismas pero los enfermos tienen valores muy dispersos, las varianzas son distintas — hay heterocedasticidad.

    Levene responde: ¿la dispersión de esta variable es igual en ambos grupos?

    Sirve para entender cómo se comporta la enfermedad. Si los enfermos tienen una varianza mucho más alta que los sanos, confirmas que la retinopatía afecta a cada persona de forma muy distinta (mucha variabilidad), mientras que los sanos son un grupo muy uniforme.

    \begin{itemize}
        \item Si p-valor $< 0.05$ $\rightarrow$ Las varianzas son distintas.
        \item Si p-valor $> 0.05$ $\rightarrow$ Las varianzas son iguales.
    \end{itemize}
\end{itemize}

